\documentclass[11pt,reqno]{amsart}

\usepackage[T1]{fontenc}
\usepackage{amsmath,amssymb,amsthm}
\usepackage[margin=1.15in]{geometry}
\usepackage[colorlinks=true,linkcolor=blue,citecolor=blue,urlcolor=blue]{hyperref}

\newtheorem{theorem}{Theorem}[section]
\newtheorem{proposition}[theorem]{Proposition}
\newtheorem{lemma}[theorem]{Lemma}
\newtheorem{corollary}[theorem]{Corollary}
\theoremstyle{definition}
\newtheorem{definition}[theorem]{Definition}
\newtheorem{example}[theorem]{Example}
\newtheorem{question}[theorem]{Question}
\newtheorem{remark}[theorem]{Remark}

\DeclareMathOperator{\supp}{supp}
\DeclareMathOperator{\spann}{span}
\newcommand{\N}{\mathbb{N}}
\newcommand{\C}{\mathbb{C}}
\newcommand{\R}{\mathbb{R}}
\newcommand{\D}{\mathbb{D}}
\newcommand{\F}{\mathcal{F}}
\newcommand{\cD}{\mathcal{D}}
\newcommand{\cI}{\mathcal{I}}
\newcommand{\cS}{\mathcal{S}}
\newcommand{\cU}{\mathcal{U}}
\newcommand{\T}{\mathfrak{T}}
\newcommand{\PP}{\mathcal{P}}
\newcommand{\HH}{\mathcal{H}}
\newcommand{\pluf}{\pi}

\title[Support families of closed subspaces of $\ell^2$]
{Support families of closed subspaces of $\ell^2$,\\
with an application to diagonalizable projection filters}

\author{David V. Feldman}
\address{Department of Mathematics and Statistics, University of New Hampshire,
Durham, NH 03824, USA}
\email{dvfinnh@gmail.com}

\subjclass[2020]{Primary 46C07, 05B35, 06C15; Secondary 30J05, 81P10, 03E05,
15A03, 46B45}

\keywords{support, closed subspace, projection lattice, maximal abelian
subalgebra, maximal filter, ultrafilter, matroid, circuit, scrawl, covering
number, representability, model space, Blaschke product}

\begin{document}
\maketitle

\begin{abstract}
For a closed subspace $M$ of $\ell^2(\N)$ let $\F(M)$ be the family of supports
of its vectors. We determine which families arise this way in finite rank and
record what survives in infinite rank. Closedness forces $\F(M)$ to be closed
under arbitrary unions, and Gaussian elimination among minimal supports is the
circuit elimination axiom; combined with a theorem of Aroca, Bossinger,
Falkensteiner, Garay L\'opez, Gonz\'alez-Ram\'irez and Valencia Negrete, this
gives, over $\C$, exactly the scrawl families of cofinitary matroids on $\N$
with $\C$-representable duals of finite rank. Of the three Bowler--Carmesin
scrawl axioms, union closure holds outright and elimination holds for finitely
many deleted coordinates; a model space $K_\Theta$ shows the maximality axiom
can fail, and can fail so far that $\F(K_\Theta)$ is the scrawl family of no
matroid. The support family also computes the interaction of $M$ with the
diagonal masa: the trace of a maximal filter of the projection lattice on the
diagonal block is $\bigcap_{M}\cD(M)$, matroid independence is
lattice-theoretic disjointness from the block, and $M$ belongs to a maximal
filter with ultrafilter trace exactly when the covering number $\chi(M)$ of its
support matroid is infinite. This invariant grades a property studied in the
companion papers: a subspace is intimate exactly when $\chi\ge3$. The grading is
strict --- Gowers's intimate subspace has $\chi=3$, so it lies in no
diagonalizable maximal filter --- but collapses on filters, every member of a
diagonalizable maximal filter having $\chi=\infty$.
\end{abstract}

\section{Introduction}\label{sec:intro}

For $x\in\ell^2(\N)$ write $\supp(x)=\{n:x_n\neq0\}$, and for a closed subspace
$M\subseteq\ell^2(\N)$ put
\[
\F(M)=\{\supp(x):x\in M\}\subseteq2^{\N}.
\]

\begin{question}\label{q:main}
Which families $\F\subseteq2^{\N}$ arise as $\F(M)$?
\end{question}

The answer in finite rank is a piece of matroid representability
(Section~\ref{sec:finiterank}). In infinite rank it is open, and
Section~\ref{sec:infinite} locates the obstruction in a question about
coefficient vanishing in model spaces.

The reason for asking is lattice-theoretic, and it is the reason this paper
belongs with \cite{FWI,FWII,FWIII,FWIV}. Write $\PP(\ell^2)$ for the lattice of
closed subspaces, and let $\{\ell^2(S):S\subseteq\N\}$ be the projection
lattice of the diagonal masa --- a copy of $2^{\N}$, and a maximal Boolean
subalgebra, a \emph{block} in the language of quantum logic \cite{K83}. A
\emph{pluf} is a maximal filter in $\PP(\ell^2)$ \cite{FWI}; its trace on the
block is $\cS(\pluf)=\{S:\ell^2(S)\in\pluf\}$, and the question of when that
trace is an ultrafilter is the diagonalization question of \cite{FWII}.
Section~\ref{sec:lattice} shows that the trace is computed from support
families and from nothing else, that matroid independence is disjointness from
the block, and that a single invariant --- the covering number of the support
matroid --- decides whether $M$ can sit in a pluf with ultrafilter trace.

Section~\ref{sec:series} draws the consequences for the companion papers. The
covering number grades the intimacy of \cite{FWII}: a subspace is intimate
exactly when $\N$ is not the union of two sets $S$ with $M\wedge\ell^2(S)=0$,
that is, exactly when $\chi(M)\ge3$, while membership in a diagonalizable pluf
requires $\chi(M)=\infty$. The two conditions differ on subspaces:
Gowers's intimate subspace of \cite{FWII} has $\chi=3$
(Theorem~\ref{thm:gowerschi}), so it belongs to no diagonalizable pluf. They
agree on plufs: every member of a pluf with ultrafilter trace has
$\chi=\infty$ (Proposition~\ref{prop:collapse}), so ``every member intimate''
and ``every member diagonally consistent'' define the same class of filters.
Section~\ref{sec:series} closes with a constraint on the rigidity conjecture of
\cite{FWII}: every addable blocker has infinite covering number.

\subsection*{Conventions and the choice of field}
$\ell^2=\ell^2(\N)$ is taken over $\C$. Two of the results below are
field-sensitive and one is not; we fix which is which.
Union closure (Section~\ref{sec:unions}) needs only that there are more scalars
than coordinates, and holds over $\R$ and $\C$ alike; over a finite field it
fails, and Section~\ref{sec:finiterank} shows the failure is not an artifact of
the argument. The finite-rank classification is representability over the field
in question, and representability genuinely depends on it
(Question~\ref{q:realcomplex}); the theorem of \cite{ABFGGV} that we invoke
carries a cardinality hypothesis on the field which $\C$ satisfies. Everything
in Sections~\ref{sec:lattice}--\ref{sec:series} --- the Galois connection, the
trace formula, the covering criterion, and the consequences for
\cite{FWII} --- uses no property of the scalars at all and holds verbatim over
$\R$. The companion papers work over $\R$, where their formal verification is
carried out; the results transferred here are among those to which the field is
immaterial.

Members of $\F(M)$ include $\emptyset=\supp(0)$. For $S\subseteq\N$ we write
$\ell^2(S)=\{x:\supp(x)\subseteq S\}$, which is $\HH_S$ in the notation of
\cite{FWII}.

\section{Closure under unions}\label{sec:unions}

\begin{lemma}\label{lem:finiteunion}
Let $M$ be any linear subspace of $\ell^2$ and $x,y\in M$. Then
$\supp(x)\cup\supp(y)\in\F(M)$.
\end{lemma}

\begin{proof}
For $\lambda\neq0$ the coordinate $(x+\lambda y)_n$ vanishes for at most one
$\lambda$, namely $-x_n/y_n$ when $y_n\neq0$. There are countably many
coordinates and uncountably many scalars, so some $\lambda$ avoids every bad
value, and then $\supp(x+\lambda y)=\supp(x)\cup\supp(y)$.
\end{proof}

The proof uses that there are more scalars than coordinates. The infinite
version uses closedness, and uses it once.

\begin{proposition}[Countable unions]\label{prop:countableunion}
Let $M\subseteq\ell^2$ be closed and $x^{(1)},x^{(2)},\dots\in M$. Then there
is $x\in M$ with $\supp(x)=\bigcup_i\supp(x^{(i)})$.
\end{proposition}

\begin{proof}
Normalize $\|x^{(i)}\|_2=1$, so $\|x^{(i)}\|_\infty\le1$. We choose $c_i\neq0$
recursively and set $x=\sum_ic_ix^{(i)}$, writing $s_m=\sum_{i\le m}c_ix^{(i)}$.

At stage $m$, two constraints bind. First, to keep
$\supp(s_m)=\supp(s_{m-1})\cup\supp(x^{(m)})$, the coefficient $c_m$ must avoid
$-(s_{m-1})_n/x^{(m)}_n$ for each of the countably many $n$ with
$x^{(m)}_n\neq0$. Second, having chosen $c_m$, put
\[
\delta_m=\min\{|(s_m)_n|:n\le m,\ (s_m)_n\neq0\}>0,
\]
a minimum over a finite set, and require of all later coefficients that
$|c_i|\le2^{-(i-m)-1}\delta_m$ for $i>m$, together with $|c_i|\le2^{-i}$
throughout. The admissible set for $c_m$ is a punctured disc less a countable
set, so the recursion runs.

Since $\sum_i|c_i|\,\|x^{(i)}\|_2<\infty$ the series converges in $\ell^2$; the
partial sums lie in $M$ and $M$ is closed, so $x\in M$. This is the only step
that uses closedness. Every term is supported in $\bigcup_i\supp(x^{(i)})$, so
$\supp(x)$ is contained in that union. Conversely, if $n\in\supp(x^{(m)})$ then
$n\in\supp(s_m)$, so $|(s_m)_n|\ge\delta_{\max(m,n)}$ once the index is large
enough for $n$ to be frozen, and the tail $\sum_{i>\max(m,n)}|c_i|$ is at most
half of it; hence $x_n\neq0$.
\end{proof}

\begin{corollary}\label{cor:arbunion}
$\F(M)$ is closed under arbitrary unions, and $\sigma(M):=\bigcup\F(M)$ is
itself a member of $\F(M)$.
\end{corollary}

\begin{proof}
An arbitrary union of supports is a union of countably many of them: a subset
of $\N$ that is a union of members of $\F(M)$ is the union of countably many,
one through each of its points. Apply
Proposition~\ref{prop:countableunion}.
\end{proof}

Closedness is not decoration here. The subspace of finitely supported sequences
orthogonal to a fixed vector of full support is a non-closed subspace whose
support family omits its own union.

\section{Elimination}\label{sec:elimination}

Nonemptiness and union closure do not suffice.

\begin{example}\label{ex:elimination}
$\F=\{\emptyset,\{1,2\},\{1,3\},\{1,2,3\}\}$ contains $\emptyset$ and is closed
under unions, but is no $\F(M)$: vectors $x,y$ with supports $\{1,2\}$ and
$\{1,3\}$ have $x_1,y_1\neq0$, and $z=x-(x_1/y_1)y$ has support $\{2,3\}$.
\end{example}

\begin{proposition}[Elimination]\label{prop:elimination}
Let $A\neq B$ be nonempty members of $\F(M)$ and $p\in A\cap B$. Then
$(A\cup B)\setminus\{p\}$ contains a nonempty member of $\F(M)$. Minimality of
$A$ and $B$ is not required; the case of minimal $A,B$ is the circuit
elimination axiom.
\end{proposition}

\begin{proof}
Eliminate $p$ as in Example~\ref{ex:elimination}. The result is nonzero, since
$z=0$ would make $x$ a multiple of $y$ and force $A=B$, and is supported in
$(A\cup B)\setminus\{p\}$.
\end{proof}

Nonemptiness, union closure and elimination among minimal members are the
axioms of a structure introduced for unrelated reasons, to which we now turn.

\section{Circuits and scrawls}\label{sec:matroids}

We use only the following. A \emph{matroid} on a ground set $E$ is given by its
family of \emph{circuits}: a nonempty antichain $\mathcal C\subseteq2^E$
satisfying circuit elimination --- for distinct $C,C'\in\mathcal C$ and
$p\in C\cap C'$, some circuit lies in $(C\cup C')\setminus\{p\}$ --- together
with the requirement that the \emph{independent} sets, those containing no
circuit, satisfy the maximality axiom (SM): every independent subset of a set
$X$ extends to a maximal independent subset of $X$. For finite $E$ the last
condition is automatic; for infinite $E$ it is the axiom isolated by Bruhn,
Diestel, Kriesell, Pendavingh and Wollan \cite{BDKPW13}. A \emph{scrawl} is a
union of circuits, and the scrawl family is closed under arbitrary unions.
Bowler and Carmesin \cite{BC18} axiomatize matroids by their scrawls: (S1)
union closure; (S2) scrawl elimination, the multi-coordinate form of circuit
elimination; and (SM).

For a matrix $A$ over a field $K$ with columns indexed by $E$, the supports of
the vectors of $\ker A$ are the scrawls of the matroid whose circuits are the
minimal dependencies among the columns. Question~\ref{q:main} asks which scrawl
families are realized when $K=\C$, $E=\N$, and the kernel is required to be
norm-closed in $\ell^2$.

\begin{theorem}\label{thm:S1S2}
Let $M\subseteq\ell^2(\N)$ be closed. Then $\F(M)$ satisfies \textup{(S1)}, and
satisfies \textup{(S2)} for finite deleted sets: if $w=\supp(x)$ with $x\in M$,
if $X\subseteq w$ is finite, and if $u_p\in M$ for $p\in X$ have supports $w_p$
with $p\in w_q\iff p=q$, then for every
$z\in w\setminus\bigcup_{p\in X}w_p$ there is $y\in M$ with
\[
z\in\supp(y)\subseteq\Bigl(w\cup\bigcup_{p\in X}w_p\Bigr)\setminus X.
\]
\end{theorem}

\begin{proof}
(S1) is Corollary~\ref{cor:arbunion}. For the elimination, $a_p=x_p/(u_p)_p$ is
well defined by the hypothesis on the $w_q$, and
$y=x-\sum_{p\in X}a_pu_p$ --- a finite sum, hence in $M$ --- has $y_p=0$ for
$p\in X$. No subtracted term touches $z$, so $y_z=x_z\neq0$.
\end{proof}

For infinite $X$ the correction term is an infinite series that must converge
in $\ell^2$, and there is no freedom to arrange it: rescaling $u_p$ rescales
$a_p$ inversely, so $a_pu_p$ is determined by the data.
Question~\ref{q:S2infinite} records the gap. Axiom (SM) is the subject of
Section~\ref{sec:infinite}.

\section{The finite-rank answer}\label{sec:finiterank}

\begin{theorem}[Aroca et al.\ {\cite[Thms.~2.15, 3.5]{ABFGGV}}]\label{thm:ABFGGV}
Let $K$ be a field, $E$ a set, and $\{0\}\neq W\subseteq K^E$ finite
dimensional. Write $\T(W)$ for the supports of elements of $W$. Then the
minimal nonempty members of $\T(W)$ are the circuits of a matroid on $E$ and
every member of $\T(W)$ is a union of circuits; and if $\#E<\#K$ then $\T(W)$
is exactly the scrawl family of that matroid.
\end{theorem}

The cardinality hypothesis is sharp. Over $\mathbb F_2$, the span of $(0,1,1)$
and $(1,0,1)$ in $\mathbb F_2^3$ has support family
$\{\emptyset,\{2,3\},\{1,3\},\{1,2\}\}$, which is not union-closed: there are
too few scalars to run Lemma~\ref{lem:finiteunion}. For $E=\N$ and $K=\C$ the
hypothesis reads $\aleph_0<2^{\aleph_0}$.

\begin{theorem}\label{thm:finiterank}
Let $M\subseteq\ell^2(\N)$ be a nonzero closed subspace of finite codimension
in some $\ell^2(S)$, or more generally a closed subspace whose orthogonal
complement within its own support block is finite dimensional. Then $\F(M)$ is
the scrawl family of a matroid on $\N$ whose dual is $\C$-representable of
finite rank; conversely every such scrawl family arises. In particular every
finite restriction of every support family of a closed subspace is the scrawl
family of a representable matroid.
\end{theorem}

\begin{proof}
A closed subspace of finite codimension in $\ell^2(S)$ is the kernel of a
finite matrix over $\C$ acting on $\ell^2(S)$, so Theorem~\ref{thm:ABFGGV}
applies with $E=S$, $K=\C$; extend by $\emptyset$ off $S$. Conversely a
representable dual of finite rank presents $M$ as such a kernel, and a kernel
of finitely many continuous functionals is closed. For the last sentence,
restrict to a finite set of coordinates: the restriction of $\F(M)$ is the
support family of the image of $M$ under the coordinate projection, a
finite-dimensional space over $\C$.
\end{proof}

By Mayhew, Newman and Whittle \cite{MNW14}, following V\'amos \cite{V78}, the
class of representable matroids is not finitely axiomatizable, so the
characterization in Theorem~\ref{thm:finiterank} is decidable
\cite{S87} but admits no finite list of axioms.

\section{Infinite rank: a model space}\label{sec:infinite}

\begin{lemma}\label{lem:expzeros}
Let $\mu_1<\dots<\mu_m$ be real and let $a\in\R^m$ be nonzero. Then
$F(t)=\sum_{k\le m}a_ke^{\mu_kt}$ has at most $m-1$ real zeros.
\end{lemma}

\begin{proof}
Induct on $m$, discarding vanishing coefficients so that all $a_k\neq0$. For
$m=1$, $F$ has no zeros. For $m>1$, suppose $F$ had $m$ zeros. Then
$G(t)=e^{-\mu_1t}F(t)$ has the same $m$ zeros, so by Rolle's theorem $G'$ has
at least $m-1$; but
$G'(t)=\sum_{k\ge2}a_k(\mu_k-\mu_1)e^{(\mu_k-\mu_1)t}$ is an exponential sum
with $m-1$ nonzero coefficients and distinct exponents, hence has at most
$m-2$ zeros.
\end{proof}

\begin{example}\label{ex:modelspace}
Choose distinct $\lambda_1,\lambda_2,\dots\in(0,1)$ with
$\sum_k(1-\lambda_k)<\infty$ --- the summability is used only for the
identification below, not for the support claim --- put $v^{(k)}=(\lambda_k^{\,n})_{n\ge0}$ and
$M=\overline{\spann}\{v^{(k)}\}$. Then every cofinite subset of $\N_0$ lies in
$\F(M)$.
\end{example}

\begin{proof}
Fix finite $S\subseteq\N_0$ and put $m=|S|+1$. Applying
Lemma~\ref{lem:expzeros} with $\mu_k=\log\lambda_k$ shows that a nontrivial
combination of $v^{(1)},\dots,v^{(m)}$ vanishes at no more than $m-1$ points of
$\N_0$; in particular the $v^{(k)}$ are linearly independent, so their span
$V_m$ has dimension $m$. Evaluation at the $|S|=m-1$ points of $S$ is a linear
map $V_m\to\R^{m-1}$, so some nonzero $x\in V_m$ vanishes on $S$. That $x$ has
at least $m-1$ zeros on $\N_0$ and at most $m-1$, so its zero set is exactly
$S$ and $\supp(x)=\N_0\setminus S$.
\end{proof}

Under $\ell^2(\N_0)\cong H^2(\D)$ the vector $v^{(k)}$ is the reproducing
kernel at $\lambda_k$ and $M$ is the model space $K_\Theta=H^2\ominus\Theta H^2$
for the Blaschke product $\Theta$ with zeros $\lambda_k$; members of $\F(M)$ are
Taylor-coefficient supports of functions in $K_\Theta$ \cite{GMR16}.

What Example~\ref{ex:modelspace} establishes is that no cofinite set is a
minimal nonempty member of $\F(M)$, since every cofinite set properly contains
another. Whether $\F(M)$ contains any set of infinite complement, and hence
whether it has any minimal nonempty member at all, is
Question~\ref{q:coeff}. We avoid the word \emph{circuit} for the minimal
members here: that terminology presupposes that $\F(M)$ is a scrawl family,
which is what is in doubt.

Suppose $K_\Theta$ contains no nonzero function with infinitely many vanishing
Taylor coefficients. Then every nonzero support is cofinite and
$\F(M)=\{\emptyset\}\cup\{\text{cofinite sets}\}$, and $\F(M)$ is the scrawl
family of no matroid --- for a reason weaker than the failure of (SM). Were
$\F(M)$ the family of unions of a circuit family $\mathcal C$, any $c\in\mathcal
C$ would be a nonempty member of $\F(M)$, hence cofinite; it properly contains
a cofinite set $c'$, itself a nonempty member of $\F(M)$ and so a union of a
nonempty subfamily of $\mathcal C$, which supplies some $c''\in\mathcal C$ with
$c''\subseteq c'\subsetneq c$ --- against the antichain axiom. Neither
elimination nor (SM) is used. Separately, the sets containing no nonempty
support are exactly the coinfinite ones, which have no maximal element, so
(SM) fails as well, and there is no minimal nonempty member; (S1) and (S2)
hold vacuously. This chain is conditional on Question~\ref{q:coeff};
Example~\ref{ex:modelspace} is not.

\section{The diagonal block}\label{sec:lattice}

\begin{definition}\label{def:DI}
For $M\in\PP(\ell^2)$ put
\[
\cD(M)=\{S\subseteq\N:M\wedge\ell^2(S)\neq0\},\qquad
\cI(M)=2^{\N}\setminus\cD(M).
\]
\end{definition}

\begin{lemma}\label{lem:meetnonzero}
$M\wedge\ell^2(S)\neq0$ if and only if $S$ contains a nonempty member of
$\F(M)$. Hence $\cI(M)$ consists exactly of the sets containing no nonempty
support, and $\cD(M)$ is determined by $\F(M)$.
\end{lemma}

\begin{proof}
The meet of closed subspaces is the intersection, and a nonzero
$x\in M\cap\ell^2(S)$ is a nonzero $x\in M$ with $\supp(x)\subseteq S$.
\end{proof}

In finite rank, where Theorem~\ref{thm:ABFGGV} applies, $\cI(M)$ is the family
of independent sets of the support matroid: matroid independence is
lattice-theoretic disjointness from the block. The minimal members of $\cD(M)$
are its circuits, and axiom (SM) asserts that a subspace of the block disjoint
from $M$ extends to a maximal such subspace.

\begin{proposition}\label{prop:trace}
For a filter $F$ on $\PP(\ell^2)$ with trace
$\cS(F)=\{S:\ell^2(S)\in F\}$,
\[
\cS(F)\subseteq\bigcap_{M\in F}\cD(M),
\]
with equality when $F$ is maximal.
\end{proposition}

\begin{proof}
If $S\in\cS(F)$ and $M\in F$ then $M\wedge\ell^2(S)\in F$, which is nonzero by
properness, so $S\in\cD(M)$. For the converse under maximality, suppose
$S\in\cD(M)$ for every $M\in F$. Then $\ell^2(S)$ meets every member of $F$
nontrivially, so by the maximality criterion (\cite{FWI}, Lemma 2.1) it belongs
to $F$.
\end{proof}

\begin{definition}\label{def:consistent}
$M$ is \emph{diagonally consistent} if some ultrafilter $\cU$ on $\N$ satisfies
$\cU\subseteq\cD(M)$. The \emph{covering number} $\chi(M)$ is the least $k$
such that $\N$ is the union of $k$ members of $\cI(M)$, and $\infty$ if there
is no such $k$.
\end{definition}

\begin{proposition}\label{prop:consistentiff}
$M$ is diagonally consistent if and only if there is a pluf $\pluf$ with
$M\in\pluf$ and $\cS(\pluf)$ an ultrafilter.
\end{proposition}

\begin{proof}
Given $\cU\subseteq\cD(M)$, let $G$ be the upward closure of the finite meets
of $\{M\}\cup\{\ell^2(S):S\in\cU\}$. Such a meet is
$M\wedge\ell^2(S_1\cap\dots\cap S_k)$, nonzero because the intersection lies in
$\cU\subseteq\cD(M)$; so $G$ is proper, and any pluf $\pluf\supseteq G$ has
$\cU\subseteq\cS(\pluf)$. A trace contains no complementary pair, since
$\ell^2(S)\wedge\ell^2(\N\setminus S)=0$, so $\cS(\pluf)=\cU$. Conversely
Proposition~\ref{prop:trace} gives $\cS(\pluf)\subseteq\cD(M)$ for $M\in\pluf$.
\end{proof}

\begin{theorem}\label{thm:covering}
$M$ is diagonally consistent if and only if $\chi(M)=\infty$.
\end{theorem}

\begin{proof}
If $\chi(M)=\infty$, the ideal $J$ generated by the downward closed family
$\cI(M)$ consists of the finite unions of its members and is proper, so its
dual filter $J^*$ is proper; extend $J^*$ to an ultrafilter $\cU$. For
$S\in\cI(M)$ we have $\N\setminus S\in J^*\subseteq\cU$, so $S\notin\cU$; hence
$\cU\subseteq\cD(M)$. Conversely, if $\cU\subseteq\cD(M)$ and
$\N=S_1\cup\dots\cup S_k$ with every $S_i\in\cI(M)$, primeness of $\cU$ puts
some $S_i\in\cU\subseteq\cD(M)$, a contradiction.
\end{proof}

\begin{proposition}\label{prop:partitionregular}
$M$ is diagonally consistent if and only if $\F(M)\setminus\{\emptyset\}$ is
partition regular: every partition of $\N$ into finitely many cells has a cell
containing a nonempty support.
\end{proposition}

\begin{proof}
Theorem~\ref{thm:covering} through Lemma~\ref{lem:meetnonzero}, covers being
refinable to partitions.
\end{proof}

For finite ground sets Edmonds' theorem \cite{E65} computes the covering number
as $\max_{A\neq\emptyset}\lceil|A|/r(A)\rceil$; see \cite[\S11.3]{Ox11}.

\section{The covering number in the pluf series}\label{sec:series}

Fix the standard basis. A closed subspace $W$ is \emph{intimate} \cite{FWII} if
for every $A\subseteq\N$ at least one of $W\wedge\ell^2(A)$,
$W\wedge\ell^2(\N\setminus A)$ is nonzero, and a nonprincipal pluf is
\emph{diagonalizable} if its trace is an ultrafilter.

\begin{proposition}\label{prop:intimatechi}
A nonzero closed subspace $W$ is intimate if and only if $\chi(W)\ge3$.
\end{proposition}

\begin{proof}
$W$ fails to be intimate exactly when some $A$ has both $A$ and
$\N\setminus A$ in $\cI(W)$, that is, exactly when $\N$ is the union of two
members of $\cI(W)$. Since $\cI(W)$ is downward closed, a cover by two members
may be replaced by a complementary pair, so this is $\chi(W)\le2$; and
$\chi(W)\ge2$ because $W\neq0$ puts $\N$ in $\cD(W)$.
\end{proof}

Intimacy is thus the first nontrivial value of an invariant whose extreme value
is diagonal consistency. The two conditions are different, and
Theorem~\ref{thm:gowerschi} separates them; but on plufs they coincide.

\begin{proposition}\label{prop:collapse}
Let $\pluf$ be a pluf whose trace is an ultrafilter. Then every member of
$\pluf$ is diagonally consistent, hence has $\chi=\infty$.
\end{proposition}

\begin{proof}
Proposition~\ref{prop:trace} gives $\cS(\pluf)\subseteq\cD(M)$ for each
$M\in\pluf$, and $\cS(\pluf)$ is an ultrafilter; apply
Theorem~\ref{thm:covering}.
\end{proof}

\begin{corollary}\label{cor:collapse}
For a nonprincipal pluf the following are equivalent: it is diagonalizable;
every member is intimate; every member is diagonally consistent.
\end{corollary}

\begin{proof}
The first two are equivalent by \cite{FWII}, Theorem 5.2. The first implies the
third by Proposition~\ref{prop:collapse}, and the third implies the second by
Proposition~\ref{prop:intimatechi}.
\end{proof}

The equivalence is a property of filters, not of subspaces. The subspace
constructed by Gowers in \cite{FWII} is intimate and has the least covering
number compatible with that.

\begin{theorem}\label{thm:gowerschi}
Let $X$ be the intimate subspace of \cite{FWII}, Theorem 4.2. Then
$\chi(X)=3$. Consequently $X$ belongs to no diagonalizable pluf.
\end{theorem}

\begin{proof}
Write the coordinates in pairs $P_i=\{2i-1,2i\}$, $i\ge1$. By \cite{FWII},
Remark 4.4 the membership condition for $X$ is that the sequence of pair sums
$s_i=x_{2i-1}+x_{2i}$ be orthogonal to every $u_n$, and by \cite{FWII},
Remark 4.7 the orthogonal complement of $\{u_n\}$ is the line spanned by
$w=(1,\tfrac12,\tfrac13,\dots)$. So $x\in X$ if and only if $s=cw$ for some
scalar $c$.

We first identify $\cI(X)$. If $S\supseteq P_i$ for some $i$, then
$e_{2i-1}-e_{2i}\in X\wedge\ell^2(S)$, so $S\in\cD(X)$. If $S$ meets every
$P_i$ in exactly one point $a_i$, the vector with $x_{a_i}=1/i$ has pair sums
$s_i=1/i$, hence lies in $X\wedge\ell^2(S)$, so again $S\in\cD(X)$. If $S$
contains no $P_i$ and misses some $P_j$ entirely, then any
$x\in X\wedge\ell^2(S)$ has $s_j=0$, forcing $c=0$ and $s\equiv0$; then for
each $i$ the two coordinates of $P_i$ sum to zero while at most one of them is
nonzero, so both vanish, and $x=0$. Hence
\[
\cI(X)=\{S:\ S\text{ contains no }P_i\text{ and misses some }P_j\}.
\]

Now $\chi(X)\ge3$ by Proposition~\ref{prop:intimatechi}, $X$ being intimate.
For the upper bound, take
\[
S=\{1,6\}\cup\{n\ \text{odd}:n\ge7\},\quad
T=\{2,3\}\cup\{n\ \text{even}:n\ge8\},\quad
V=\{4,5\}.
\]
These cover $\N$. None contains a complete pair --- each meets $P_i$ in at most
one point --- and each misses one entirely: $S$ misses $P_2$, $T$ misses $P_3$,
and $V$ misses $P_1$. So all three lie in $\cI(X)$ and $\chi(X)\le3$. The last
claim is Proposition~\ref{prop:collapse}.
\end{proof}

\begin{proposition}\label{prop:blocker}
Let $\cU$ be a non-selective ultrafilter on $\N$ and let $R$ be an addable
blocker in the sense of \cite{FWII}, \S3: a closed subspace with
$\dim(R\wedge\ell^2(S))=\infty$ for every $S\in\cU$. Then $\chi(R)=\infty$.
\end{proposition}

\begin{proof}
Each $S\in\cU$ has $R\wedge\ell^2(S)\neq0$, so $\cU\subseteq\cD(R)$; apply
Theorem~\ref{thm:covering}.
\end{proof}

This is a fourth constraint on the rigidity conjecture of \cite{FWII},
independent of the three proved in \cite{FWIV}, \S6: those bound the dimension
of a blocker on finitely many pieces, place its total support in $\cU$, and
force infinite piece-spread inside every measure-one block, while
Proposition~\ref{prop:blocker} forbids finite covers of $\N$ by sets disjoint
from $R$.

\section{Questions}\label{sec:questions}

\begin{question}\label{q:mainopen}
Characterize the support families of closed subspaces of $\ell^2$ of infinite
rank.
\end{question}

\begin{question}\label{q:S2infinite}
Does $\F(M)$ satisfy scrawl elimination \textup{(S2)} when the deleted set is
infinite?
\end{question}

\begin{question}\label{q:coeff}
For a Blaschke product $\Theta$, which subsets of $\N_0$ occur as
$\{n:\hat f(n)\neq0\}$ for $f\in K_\Theta$? Does $K_\Theta$ contain a nonzero
function with infinitely many vanishing Taylor coefficients? By
Proposition~\ref{prop:partitionregular}, $K_\Theta$ is diagonally consistent
exactly when its family of nonzero coefficient supports is partition regular,
so this question decides whether $K_\Theta$ can sit in a pluf with ultrafilter
trace.
\end{question}

\begin{question}\label{q:spectrum}
Which values does $\chi(M)$ take, and how is it computed? Edmonds' formula
governs finite ground sets; Theorem~\ref{thm:gowerschi} exhibits the value $3$
on an infinite-rank example.
\end{question}

\begin{question}\label{q:traces}
Which filters on $\N$ arise as traces of plufs? By \cite{FWII}, Theorem 3.4 the
block filter of $\cU$ is itself maximal exactly for selective $\cU$; the
present question asks for the traces of arbitrary plufs.
\end{question}

\begin{question}\label{q:realcomplex}
Theorem~\ref{thm:finiterank} is a statement about $\C$-representability. Which
support families of closed subspaces of the real $\ell^2$ arise, and where do
the two classifications differ?
\end{question}

\begin{question}\label{q:otherblocks}
Replace the diagonal masa by another maximal Boolean subalgebra of
$\PP(\ell^2)$. For the continuous masa the analogous question is the one posed
in \cite{FWII}, \S6.
\end{question}

\subsection*{Formal verification}
Sections~\ref{sec:unions}, \ref{sec:elimination}, \ref{sec:lattice} and
\ref{sec:series} have been formalized and machine-checked in Lean~4 with
Mathlib (toolchain v4.28.0), extending the development that accompanies
\cite{FWI,FWII,FWIII,FWIV}: Lemma~\ref{lem:finiteunion},
Proposition~\ref{prop:countableunion} with its recursion on frozen minima,
Corollary~\ref{cor:arbunion} and Proposition~\ref{prop:elimination}, whose
independence of minimality the mechanization established;
Lemma~\ref{lem:meetnonzero}, both halves of Proposition~\ref{prop:trace},
Proposition~\ref{prop:consistentiff} and Theorem~\ref{thm:covering}; and all
of Section~\ref{sec:series} --- Propositions~\ref{prop:intimatechi},
\ref{prop:collapse} and \ref{prop:blocker}, Corollary~\ref{cor:collapse}, and
Theorem~\ref{thm:gowerschi} with the three-cover displayed above. The
artifacts compile with no unproved obligations, and every statement's axiom
footprint is exactly \texttt{propext}, \texttt{Classical.choice},
\texttt{Quot.sound}.

The mechanization is carried out over $\R$, which the conventions of
Section~\ref{sec:intro} license: no item required a property of the scalars
beyond uncountability of the field, used once, in the counting step behind
Lemma~\ref{lem:finiteunion} and Proposition~\ref{prop:countableunion}.

Sections~\ref{sec:matroids}--\ref{sec:infinite} are covered as well.
Mathlib's matroid API has no constructor building a matroid from a circuit
family, and no scrawls, cofinitary duals or representability; the
mechanization therefore defines circuit families, independence, the maximality
axiom and scrawl families directly, in terms of families of subsets of $\N$,
and imports nothing from that API. Against those definitions:
Theorem~\ref{thm:S1S2} in both clauses; Theorem~\ref{thm:finiterank},
conditional on Theorem~\ref{thm:ABFGGV}, which is supplied as a named
hypothesis about the support families of finite-dimensional subspaces of
$\R^{\N}$ and never discharged --- the transfer between those and the support
families of finite-rank closed subspaces of $\ell^2$, in both directions, is
proved, the realization being by a diagonal rescaling with nonvanishing
entries; Lemma~\ref{lem:expzeros}, for which a counting form of Rolle's
theorem was built, Mathlib having none; and Example~\ref{ex:modelspace} and
the conditional consequence drawn from it, in the strengthened form recorded
in Section~\ref{sec:infinite}: the mechanized argument uses only the antichain
axiom.

What remains outside the development is what should: Theorem~\ref{thm:ABFGGV},
which is a theorem of the literature, and Question~\ref{q:coeff}, which is
open.

\subsection*{Acknowledgements}
Portions of these results were developed in the course of an extended dialogue
with Anthropic's Claude language model; responsibility for correctness rests
with the author.

\begin{thebibliography}{ABFGGV}

\bibitem[AB15]{AB15}
S.~H. Afzali Borujeni and N.~Bowler,
\emph{Thin sums matroids and duality},
Adv. Math. \textbf{271} (2015), 1--29.

\bibitem[ABFGGV]{ABFGGV}
F.~Aroca, L.~Bossinger, S.~Falkensteiner, C.~Garay L\'opez,
L.~R. Gonz\'alez-Ram\'irez and C.~V. Valencia Negrete,
\emph{Infinite matroids in tropical differential algebra},
Canad. Math. Bull.; \texttt{arXiv:2305.04784}.

\bibitem[BC18]{BC18}
N.~Bowler and J.~Carmesin,
\emph{An excluded minors method for infinite matroids},
J. Combin. Theory Ser. B \textbf{128} (2018), 104--113.

\bibitem[BDKPW13]{BDKPW13}
H.~Bruhn, R.~Diestel, M.~Kriesell, R.~Pendavingh and P.~Wollan,
\emph{Axioms for infinite matroids},
Adv. Math. \textbf{239} (2013), 18--46.

\bibitem[E65]{E65}
J.~Edmonds,
\emph{Minimum partition of a matroid into independent subsets},
J. Res. Nat. Bur. Standards Sect.~B \textbf{69B} (1965), 67--72.

\bibitem[FWI]{FWI}
D.~V. Feldman and A.~Wilce, \emph{Ultrafilter limits in the projection lattice
of a Hilbert space}, preprint (paper I of this series).

\bibitem[FWII]{FWII}
D.~V. Feldman and A.~Wilce, \emph{Intimate subspaces, block filters, and the
diagonalization of maximal projection filters}, preprint (paper II).

\bibitem[FWIII]{FWIII}
D.~V. Feldman and A.~Wilce, \emph{Projection-lattice ultrafilters at a
measurable cardinal}, preprint (paper III).

\bibitem[FWIV]{FWIV}
D.~V. Feldman and A.~Wilce, \emph{Exact paving for Fubini products of normal
measures, with an application to maximal projection filters}, preprint
(paper IV).

\bibitem[GMR16]{GMR16}
S.~R. Garcia, J.~Mashreghi and W.~T. Ross,
\emph{Introduction to Model Spaces and their Operators},
Cambridge Univ. Press, 2016.

\bibitem[K83]{K83}
G.~Kalmbach, \emph{Orthomodular Lattices},
London Math. Soc. Monographs \textbf{18}, Academic Press, 1983.

\bibitem[MNW14]{MNW14}
D.~Mayhew, M.~Newman and G.~Whittle,
\emph{Is the missing axiom of matroid theory lost forever?},
Quart. J. Math. \textbf{65} (2014), 1397--1415.

\bibitem[Ox11]{Ox11}
J.~Oxley, \emph{Matroid Theory}, 2nd ed., Oxford Univ. Press, 2011.

\bibitem[S87]{S87}
B.~Sturmfels,
\emph{On the decidability of Diophantine problems in combinatorial geometry},
Bull. Amer. Math. Soc. \textbf{17} (1987), 121--124.

\bibitem[V78]{V78}
P.~V\'amos,
\emph{The missing axiom of matroid theory is lost forever},
J. London Math. Soc. \textbf{18} (1978), 403--408.

\end{thebibliography}

\end{document}
